\chapter{Heavy Tails and Extreme Values}
\label{ch:heavytails}

\begin{tcolorbox}[feynbox={Sky}{BBg},title={Before and After}]
\textbf{Before this chapter you need:} logarithms and log--log plots
(\S\ref{sec:logrules}, \S\ref{sec:logscale}), integrals
(\S\ref{sec:integrals}), moments and their existence
(\S\ref{sec:variance}), the law of large numbers and the CLT
(\S\ref{sec:lln}--\S\ref{sec:clt}), estimators and consistency
(\S\ref{sec:estimators}), joint distributions (\S\ref{sec:joint}).\\
\textbf{After it you will understand:} the single rule
$\E|X|^k < \infty \iff k < \alpha$ that organises this entire chapter, how the
Hill estimator is derived and why it survives where kurtosis does not, why one
degenerate series in this project produced $\hat\alpha = 31{,}581$, and what
copulas add for the multivariate case.
\end{tcolorbox}

Almost everything in this chapter follows from one rule about which moments
exist. Section~\ref{sec:moments} states it; the remaining sections are
consequences.

%% ─────────────────────────────────────────────────────────────────────────────
\section{Moment Existence: The Organising Principle}
\label{sec:moments}

\situation{
You are averaging the wealth of people in a room. Add a hundred ordinary people
and the average barely moves. Add one billionaire and it leaps. Add a second and
it leaps again. The average never settles, because the next arrival can always
outweigh everyone already counted.
}

\problem{
Standard statistics assumes averages settle. When the underlying distribution has
a heavy enough tail they do not --- not slowly, but never --- and every downstream
quantity built on that average inherits the failure. Knowing exactly which
averages are safe is therefore the first question, not a technicality.
}

\workedarith{
Suppose $\Prob(|X| > x) = x^{-\alpha}$ for $x \geq 1$. The density is
$f(x) = \alpha x^{-\alpha-1}$, and the $k$-th moment is
\[
  \E[|X|^k] = \int_1^{\infty} x^k \cdot \alpha x^{-\alpha-1}\,dx
            = \alpha\int_1^{\infty} x^{k-\alpha-1}\,dx .
\]

The integral $\int_1^\infty x^{p}\,dx$ converges only when $p < -1$. Here
$p = k - \alpha - 1$, so convergence requires
\[
  k - \alpha - 1 < -1 \iff k < \alpha .
\]

\medskip
\textbf{Concretely, at $\alpha = 3$:}
\[
  \E[|X|^2] = 3\int_1^\infty x^{-2}\,dx = 3\Bigl[-x^{-1}\Bigr]_1^\infty = 3 \quad\text{(finite)},
\]
\[
  \E[|X|^4] = 3\int_1^\infty x^{0}\,dx = 3\Bigl[x\Bigr]_1^\infty = \infty \quad\text{(divergent)}.
\]

\medskip
\textbf{Applied to this project.} BRICS tail indices run
$\hat\alpha \in [2.54, 3.21]$:

\medskip
\begin{tabular}{@{}llll@{}}
\toprule
$k$ & Quantity & Needs $\alpha >$ & Status across BRICS \\
\midrule
1 & mean       & 1 & finite everywhere \\
2 & variance   & 2 & finite everywhere \\
3 & skewness   & 3 & finite in 1 of 5 markets \\
4 & kurtosis   & 4 & \textbf{infinite everywhere} \\
\bottomrule
\end{tabular}

\medskip
Every remaining section of this chapter is a consequence of the last row.
}

\formaldef{
A random variable has a \textbf{power-law} (regularly varying) tail with
\textbf{tail index} $\alpha > 0$ if
\[
  \Prob(|X| > x) \sim L(x)\,x^{-\alpha} \quad \text{as } x \to \infty ,
\]
where $L$ is slowly varying. The governing rule is
\[
  \boxed{\;\E\bigl[|X|^k\bigr] < \infty \iff k < \alpha\;}
\]
Smaller $\alpha$ means a heavier tail and fewer finite moments.
}

\analogybreak{
``Infinite'' here is not a large number. It is the statement that the defining
integral diverges, so the quantity does not exist. A sample kurtosis can always
be computed --- the arithmetic never fails --- but with $\alpha < 4$ it is not an
estimate of anything, because there is nothing for it to converge to. This
distinction is the whole of \S\ref{sec:kurtosis}.
}

\whyhere{
This rule is why \texttt{tail\_index\_diff} is a ranked fidelity metric while
\texttt{kurtosis\_diff} and \texttt{skewness\_diff} were demoted to
descriptive-only. Ranking generators on an estimate of an infinite quantity is not
a strict test; it is noise wearing a decimal point.
}

%% ─────────────────────────────────────────────────────────────────────────────
\section{Power Laws}
\label{sec:powerlaws}

\situation{
City sizes. New York has 8 million people; the next few have 4M, 3M, 2M. Halve
the population threshold and roughly twice as many cities qualify. The same
doubling holds whether you start at a million or ten thousand --- the rule has no
characteristic scale.
}

\problem{
The normal distribution has exponentially decaying tails, which makes extreme
events not merely unlikely but effectively impossible. Power-law tails decay
polynomially, which is slow enough to make extremes a routine feature of a few
thousand observations rather than a once-per-civilisation event.
}

\workedarith{
NIFTY50, 20-year, $n = 4{,}954$.

The top 5\% of $|r_t|$: 248 observations, from $1.97\%$ up to $15.91\%$. The Hill
estimate on that tail is $\hat\alpha \approx 2.8$, so
$\Prob(|r| > x) \approx C x^{-2.8}$.

\medskip
\textbf{Scale-free, checked.} Doubling the threshold multiplies the exceedance
probability by $2^{-2.8} = 0.144$ --- the same factor whether you double from 2\%
to 4\% or from 4\% to 8\%. Under a Gaussian the analogous factor collapses:
\[
  \frac{\Prob(|Z|>4)}{\Prob(|Z|>2)} = \frac{6.3\times10^{-5}}{4.6\times10^{-2}} = 0.0014 ,
  \qquad
  \frac{\Prob(|Z|>8)}{\Prob(|Z|>4)} = \frac{1.2\times10^{-15}}{6.3\times10^{-5}} = 2\times10^{-11} .
\]
The Gaussian ratio falls by eight orders of magnitude between the two doublings;
the power-law ratio does not move at all. That is the difference between
exponential and polynomial decay.

\medskip
\textbf{Straight lines.} Taking logarithms of
$\Prob(|X|>x) \approx Cx^{-\alpha}$:
\[
  \ln\Prob(|X|>x) = \ln C - \alpha\ln x ,
\]
a straight line of slope $-\alpha$ on log--log axes (\S\ref{sec:logscale}).
Straightness there is the visual signature of a power law and the slope reads off
$\alpha$ directly.
}

\formaldef{
For a power-law tail with index $\alpha$, the \textbf{survival function} is
\[
  \bar{F}(x) = \Prob(|X| > x) \sim L(x)x^{-\alpha} ,
\]
and $L$ slowly varying means $L(cx)/L(x) \to 1$ for every $c > 0$ --- so $L$
contributes no scaling of its own.

Typical values: financial returns $\alpha \approx 2$--$4$; city sizes and word
frequencies $\alpha \approx 1$ (Zipf); a Gaussian has no finite $\alpha$, its tail
decaying faster than any power.
}

\analogybreak{
The tail index is an \emph{asymptotic} property, describing behaviour as
$x \to \infty$. Any finite sample --- even 20 years of daily data --- contains only
a few hundred genuinely extreme observations, and the estimator must draw a line
between the power-law tail and the ordinary body of the distribution. Where that
line falls is a judgement call, and \S\ref{sec:hill} shows how much the answer
depends on it.
}

\whyhere{
\texttt{tail\_index\_diff} $= |\hat\alpha_{\text{real}} - \hat\alpha_{\text{syn}}|$
is the primary heavy-tail metric in the fidelity family, replacing
\texttt{kurtosis\_diff}. At $n \approx 4{,}960$ the Hill estimator has a standard
deviation near $0.035$ --- precise enough to separate generators that reproduce the
tail from those that do not.
}

%% ─────────────────────────────────────────────────────────────────────────────
\section{Why Kurtosis Is Unreliable Here}
\label{sec:kurtosis}

\situation{
You want a single number for how extreme a dataset is, and kurtosis is the
obvious candidate. You compute it on 250 days of BOVESPA and get $2.53$. You
compute it on 2{,}000 days and get $11.84$. Neither is wrong, and that is the
problem.
}

\problem{
An estimator is supposed to settle as data accumulates. This one moves further
with every additional observation, so there is no sample size at which the answer
becomes trustworthy --- and \S\ref{sec:moments} already said why.
}

\workedarith{
BOVESPA, two sample sizes:

\medskip
\begin{tabular}{@{}rrl@{}}
\toprule
Sample $n$ & Sample kurtosis & Hill $\hat\alpha$ \\
\midrule
$250$    & $2.53$  & $4.83$ (sd $1.253$) \\
$2{,}000$ & $11.84$ & $3.11$ (sd $0.316$) \\
\bottomrule
\end{tabular}

\medskip
\textbf{Kurtosis multiplied by 4.7} as $n$ grew eightfold. It is diverging.

\textbf{Hill fell from $4.83$ to $3.11$ and its standard deviation shrank by a
factor of four.} It is converging, and the movement is bias disappearing
(\S\ref{sec:estimators}).

\medskip
\textbf{Why kurtosis diverges.} The estimator is
\[
  \hat\kappa_n = \frac{n^{-1}\sum_i (x_i - \bar{x})^4}{\bigl(n^{-1}\sum_i (x_i-\bar{x})^2\bigr)^2} .
\]
The numerator is a sample mean of $X^4$. The law of large numbers
(\S\ref{sec:lln}) requires $\E[X^4]$ to be finite, and at $\alpha < 4$ it is not.
So the numerator has nothing to converge to and grows with $n$; each new record
extreme raises it again.

\medskip
\textbf{Why Hill does not.} Hill averages \emph{log spacings} of order
statistics. Under a power law these are exponentially distributed
(\S\ref{sec:hill}), and the exponential distribution has finite moments of every
order. The law of large numbers applies to them regardless of $\alpha$.
}

\formaldef{
When $\alpha < 4$, the population kurtosis $\E[(X-\mu)^4]/\sigma^4$ is undefined.
The sample kurtosis $\hat\kappa_n$ does not converge; it grows without bound, its
numerator behaving like a sum in the domain of attraction of a stable law with
index $\alpha/4 < 1$.

The Hill estimator, by contrast, is consistent and asymptotically normal for any
$\alpha > 0$, provided the number of order statistics $k$ satisfies
$k \to \infty$ and $k/n \to 0$.
}

\analogybreak{
``MOEX sample kurtosis $= 65.70$'' is not a measurement of the population
kurtosis, which is infinite. It is an artefact of this particular sample, driven
by a handful of observations --- above all 24 February 2022 at $-33.3\%$. A
different 20-year window could give 30 or 150 with equal legitimacy. The Hill
estimate of $\hat\alpha \approx 2.7$ for the same series is stable across windows.
}

\whyhere{
This is why \texttt{kurtosis\_diff} and \texttt{skewness\_diff} moved from the
seven ranked fidelity metrics to the seven descriptive ones. It is also why the
length guard in \texttt{FinancialMetrics.\_\_init\_\_} warns when the compared
series differ in length: a diverging statistic makes any kurtosis comparison
between unequal-length series biased by length alone, independently of the data.
}

%% ─────────────────────────────────────────────────────────────────────────────
\section{The Hill Estimator, Derived}
\label{sec:hill}

\situation{
You want the slope of a straight line but only trust the right-hand portion of
your data. So you discard the left, keep the largest few hundred points, and fit
using those alone.
}

\problem{
The tail index is defined by behaviour as $x \to \infty$, and any real sample
stops well short. We need an estimator that uses only the largest observations,
where the power law plausibly holds, and ignores the body where it does not.
}

\workedarith{
\textbf{Step 1: the key transformation.} Suppose
$\Prob(X > x) = (x/u)^{-\alpha}$ for $x \geq u$. Let $Y = \ln(X/u)$. Then for
$y > 0$,
\[
  \Prob(Y > y) = \Prob\bigl(X > ue^{y}\bigr) = \bigl(e^{y}\bigr)^{-\alpha} = e^{-\alpha y} .
\]
So $Y$ is \textbf{exponential} with rate $\alpha$. The logarithm has converted a
power law into an exponential --- and exponentials are easy.

\medskip
\textbf{Step 2: estimate the rate.} An exponential with rate $\alpha$ has mean
$1/\alpha$. By the law of large numbers, averaging the $Y$ values estimates
$1/\alpha$, so
\[
  \hat\alpha = \frac{1}{\bar{Y}} = \frac{k}{\sum_{i=1}^{k}\ln(X_{(n-i+1)}/X_{(n-k)})} .
\]
That is the Hill estimator. It is maximum likelihood (\S\ref{sec:mle}) for the
exponential rate.

\medskip
\textbf{Step 3: a numerical check.} Take five tail observations with threshold
$u = 2$: values $2.2, 2.7, 3.3, 4.5, 8.1$.
\[
  \ln(2.2/2) = 0.0953, \quad \ln(2.7/2) = 0.3001, \quad \ln(3.3/2) = 0.5008,
\]
\[
  \ln(4.5/2) = 0.8109, \quad \ln(8.1/2) = 1.3971 .
\]
Sum $= 3.1042$, mean $= 0.6208$, so
\[
  \hat\alpha = 1/0.6208 = 1.61 .
\]

\medskip
\textbf{Why moments do not matter.} Each $\ln(X_{(i)}/u)$ is exponential, which
has finite moments of every order. Whatever $\alpha$ is --- even below 1, where
$X$ has no mean at all --- the log spacings are perfectly well behaved. This is
exactly the property kurtosis lacks.

\medskip
\textbf{The bias--variance trade-off in $k$.} Small $k$ uses only the most extreme
points, where the power law holds best: low bias, high variance. Large $k$ reaches
into the body where the power law fails: low variance, high bias. This project
uses $k = \lfloor 0.05n\rfloor$, so $248$ points at $n = 4{,}954$.
}

\formaldef{
Order the sample $X_{(1)} \leq \cdots \leq X_{(n)}$. Using the top $k$, the
\textbf{Hill estimator} is
\[
  \hat\alpha_{\text{Hill}} = \left[\frac{1}{k}\sum_{i=1}^{k}
    \ln\frac{X_{(n-i+1)}}{X_{(n-k)}}\right]^{-1} .
\]
It is consistent and asymptotically normal, with
$\sqrt{k}(\hat\alpha - \alpha) \to \N(0, \alpha^2)$, provided $k \to \infty$ and
$k/n \to 0$. The standard error is therefore approximately
$\hat\alpha/\sqrt{k}$.

At $\hat\alpha = 3.11$ and $k = 248$: $3.11/\sqrt{248} = 0.197$.
}

\analogybreak{
Hill assumes a power-law tail. Applied to data without one --- a Gaussian, say ---
it returns a number regardless, and that number is meaningless rather than
merely large. It is also sensitive to $k$: a Hill plot of $\hat\alpha$ against $k$
should show a stable plateau, and if it does not, the power-law assumption is
suspect. Reporting a single $\hat\alpha$ without inspecting that stability is
reporting one point from a curve.
}

\whyhere{
Hill supplies \texttt{tail\_index\_diff}, the primary heavy-tail fidelity metric.
Its independence from moment existence is exactly why it replaced kurtosis: it
remains a genuine estimator on data where the fourth moment does not exist. The
next section shows what happens when its assumption is violated outright.
}

%% ─────────────────────────────────────────────────────────────────────────────
\section{When Hill Fails: the 31{,}581 Case}
\label{sec:hillfail}

\situation{
A tape measure works well until you try to measure something with no length. The
instrument does not announce the mistake --- it returns a number, and the number
looks like a measurement.
}

\problem{
A collapsed generator produces a nearly constant series. It has no tail at all.
The Hill estimator, asked for the index of a tail that does not exist, returns a
finite number, and that number then enters a ranking table alongside genuine
estimates as though it were comparable.
}

\workedarith{
\textbf{What happened in this project.} During a BOVESPA smoke test, TimeGAN
collapsed and produced a series with standard deviation $0.0006$ against a real
$0.0103$. Hill returned
\[
  \hat\alpha = 31{,}581 .
\]

\medskip
\textbf{Why.} The estimator divides by the mean log spacing
\[
  \bar{Y} = \frac{1}{k}\sum_i \ln\frac{X_{(n-i+1)}}{X_{(n-k)}} .
\]
On a near-constant series every order statistic is nearly equal, so each ratio is
nearly 1 and each logarithm nearly 0.

Concretely, if the largest values are all within $0.003\%$ of the threshold, then
$\ln(1.00003) = 3\times10^{-5}$, and
\[
  \hat\alpha = \frac{1}{3\times10^{-5}} = 33{,}000 .
\]
Same order of magnitude as observed. The estimator is dividing by very nearly
zero, and small denominators produce enormous quotients.

\medskip
\textbf{Why this is dangerous rather than merely wrong.} The metric is
$|\hat\alpha_{\text{real}} - \hat\alpha_{\text{syn}}|$. With
$\hat\alpha_{\text{real}} \approx 3$, the collapsed model scores
$|3 - 31581| = 31578$ --- catastrophically bad, which is correct. But had the
collapse been milder, giving $\hat\alpha = 3.0$ on a series with no tail
whatsoever, it would have scored a \emph{perfect} $0$ on this metric while being
degenerate. A number with no information content had entered a ranking.

\medskip
\textbf{The guard.} Real equity indices sit between $2.4$ and $4.2$. The pipeline
therefore returns \texttt{NaN} whenever
\[
  \hat\alpha > 20 ,
\]
and \texttt{NaN} ranks last by \texttt{na\_option='bottom'} --- so a failed
estimate is scored as the worst outcome rather than silently excluded. It also
returns \texttt{NaN} when the threshold order statistic is non-positive or the
mean log spacing is non-positive, both of which signal the same degeneracy.
}

\formaldef{
The Hill estimator is undefined in the limit of a degenerate distribution. If
$X$ is constant, every log spacing is $0$, $\bar{Y} = 0$, and
$\hat\alpha = 1/0$ is undefined rather than infinite.

In floating-point arithmetic the spacings are small but non-zero, so the division
succeeds and returns a large finite value carrying no information. Guarding
requires an explicit range check, because no arithmetic error occurs.
}

\analogybreak{
The guard is a heuristic, not a theorem. A genuine distribution with
$\alpha = 25$ --- extremely thin-tailed but legitimate --- would be rejected. That
trade is deliberate: no equity index has such a tail, so within this domain the
false-rejection rate is effectively zero, while the false-acceptance rate without
the guard was demonstrably not.
}

\whyhere{
This is one of three degeneracy guards protecting the evaluation. The other two
are the post-generation check on mean, standard deviation and ACF(1), and the
autoencoder $R^2$ threshold. Each exists because a collapsed model can otherwise
score \emph{well} on some metric by accident, and a benchmark that can be won by
producing nothing is not a benchmark.
}

%% ─────────────────────────────────────────────────────────────────────────────
\section{Extreme Value Theory}
\label{sec:evt}

\situation{
A dam must survive the worst flood in a century, but you have only forty years of
records. The largest flood you have seen is not the largest that can happen, and
you must say something about a level you have never observed.
}

\problem{
The CLT (\S\ref{sec:clt}) describes averages, and averages are exactly what a risk
manager does not care about. There is a parallel theory for maxima, and it says
something surprisingly strong: whatever the underlying distribution, the maximum
has only three possible limiting shapes.
}

\workedarith{
\textbf{Maxima of uniform draws.} Take $n$ values uniform on $[0,1]$. The maximum
is below $x$ only if all $n$ are, so
\[
  \Prob(M_n \leq x) = x^{n} .
\]
At $n = 100$: $\Prob(M_{100} \leq 0.95) = 0.95^{100} = 0.0059$. The maximum is
almost certainly above $0.95$, and as $n$ grows it is pushed toward 1 --- a
\emph{bounded} limit.

\medskip
\textbf{Maxima of exponential draws.} With $\Prob(X > x) = e^{-x}$,
\[
  \Prob(M_n \leq x) = \bigl(1 - e^{-x}\bigr)^{n} .
\]
Here the maximum grows like $\ln n$ without bound, but only logarithmically ---
slowly.

\medskip
\textbf{Maxima of power-law draws.} With $\Prob(X > x) = x^{-\alpha}$, the maximum
grows like $n^{1/\alpha}$. At $\alpha = 3$ and $n$ multiplied by 1{,}000, the
typical maximum multiplies by $1000^{1/3} = 10$.

\medskip
Three regimes --- bounded, logarithmic, polynomial. The Fisher--Tippett theorem
says these exhaust the possibilities.

\medskip
\textbf{Peaks over threshold.} Rather than block maxima, fix a high threshold $u$
and model the exceedances $X - u$. As $u$ rises, that distribution approaches the
generalised Pareto, and its shape parameter is $\xi = 1/\alpha$. At
$\hat\alpha = 3.11$, $\xi = 0.32$ --- firmly in the heavy-tailed regime, since
$\xi > 0$.
}

\formaldef{
\textbf{Fisher--Tippett--Gnedenko.} If suitably normalised maxima converge, the
limit is a \textbf{generalised extreme value} distribution
\[
  G(x) = \exp\left\{-\left[1 + \xi\left(\frac{x-\mu}{\sigma}\right)\right]^{-1/\xi}\right\} ,
\]
with three cases: $\xi = 0$ Gumbel (exponential-like tails), $\xi > 0$ Fréchet
(power-law tails), $\xi < 0$ Weibull (bounded).

\textbf{Pickands--Balkema--de Haan.} For a high threshold $u$, exceedances follow
approximately a \textbf{generalised Pareto} distribution
\[
  H(y) = 1 - \left(1 + \frac{\xi y}{\beta}\right)^{-1/\xi} , \qquad y > 0 ,
\]
with $\xi = 1/\alpha$ for power-law tails. Financial returns sit in the Fréchet
domain.
}

\analogybreak{
EVT describes the limit as the threshold rises, and raising the threshold discards
data. There is no escaping the trade: a higher threshold gives a better
approximation on fewer points. EVT also assumes independent observations in its
basic form, which volatility clustering violates --- extremes arrive in bursts, so
the effective sample of independent extremes is smaller than the count suggests.
}

\whyhere{
The Hill estimator is the EVT tool this project actually uses, and it is the
peaks-over-threshold approach with $k = \lfloor 0.05n\rfloor$ setting the
threshold. The connection $\xi = 1/\alpha$ means that reporting $\hat\alpha$
between $2.54$ and $3.21$ is reporting $\xi$ between $0.31$ and $0.39$: the same
measurement in the vocabulary a risk manager would use.
}

%% ─────────────────────────────────────────────────────────────────────────────
\section{Copulas and Tail Dependence}
\label{sec:copulas}

\situation{
Two swimmers, each an excellent swimmer on their own. Whether they can rescue each
other depends on something their individual abilities do not capture: whether they
get into trouble at the same time.
}

\problem{
Section~\ref{sec:joint} established that marginals do not determine the joint
distribution, and \S\ref{sec:multivariate} showed the portfolio consequences. What
was missing is a device that separates the two questions cleanly --- and,
critically, one that can express dependence concentrated in the tails, which
correlation cannot.
}

\workedarith{
\textbf{Correlation misses the thing that matters.} Consider two return series,
each standard normal, in two scenarios:
\begin{itemize}
  \item \emph{Gaussian dependence} with correlation $0.5$;
  \item \emph{$t$-copula dependence}, also correlation $0.5$, with 3 degrees of
        freedom.
\end{itemize}
Identical marginals, identical correlation. But the probability that both fall
below their 1\% quantile simultaneously differs sharply: for the Gaussian case it
tends to zero as the quantile becomes more extreme, while for the $t$-copula it
tends to a positive constant. Joint crashes are asymptotically impossible under
one and routine under the other, with every number a correlation matrix records
being identical.

\medskip
\textbf{The transformation.} Apply each variable's own CDF:
$U = F_X(X)$ and $V = F_Y(Y)$. Whatever the original distributions, $U$ and $V$
are uniform on $[0,1]$ --- so all marginal information has been stripped out and
what remains in the joint distribution of $(U,V)$ is pure dependence. That joint
distribution is the copula.

\medskip
\textbf{Tail dependence, defined by a limit.} The lower tail dependence
coefficient is
\[
  \lambda_L = \lim_{q \to 0^{+}} \Prob\bigl(V \leq q \mid U \leq q\bigr) .
\]
The Gaussian copula has $\lambda_L = 0$ for every correlation below 1. The
$t$-copula has $\lambda_L > 0$. This single number is what the 2008 credit crisis
literature identified as missing from Gaussian-copula pricing models.
}

\formaldef{
\textbf{Sklar's theorem.} Any joint CDF $F$ with continuous marginals
$F_1,\ldots,F_d$ can be written uniquely as
\[
  F(x_1,\ldots,x_d) = C\bigl(F_1(x_1),\ldots,F_d(x_d)\bigr) ,
\]
where the \textbf{copula} $C$ is a joint CDF on $[0,1]^d$ with uniform marginals.
Marginals and dependence are thus fully separable: choose each independently.

\textbf{Tail dependence:}
\[
  \lambda_U = \lim_{q\to1^{-}}\Prob\bigl(V > q \mid U > q\bigr) ,
  \qquad
  \lambda_L = \lim_{q\to0^{+}}\Prob\bigl(V \leq q \mid U \leq q\bigr) .
\]
}

\analogybreak{
Copulas separate marginals from dependence, which is their strength and their
trap. The separation is a mathematical fact, not a licence to pick the two
independently and assume the result is realistic --- and estimating tail
dependence requires observations in the joint tail, which are by definition the
scarcest data available. A copula fitted mostly on the body of the distribution
tells you very little about $\lambda_L$.
}

\whyhere{
This section is scope-setting rather than method. This thesis is univariate: one
model per market, evaluated per market, with no cross-market dependence modelled
or measured. Copulas are the natural device for the multivariate extension, and
tail dependence would be the metric that matters --- since the question a portfolio
owner asks is whether the markets crash together, which nothing in the current
evaluation addresses.
}

%% ─────────────────────────────────────────────────────────────────────────────
\section{Wasserstein Distance}
\label{sec:wasserstein}

\situation{
Two piles of sand on a table. One is the distribution of real returns, the other
of synthetic returns. How different are they? The Wasserstein distance asks for
the minimum work --- mass times distance --- needed to reshape one pile into the
other.
}

\problem{
The Kolmogorov--Smirnov statistic measures the largest gap between two CDFs and
was designed for independent data. Market returns are not independent: the ACF of
$|r_t|$ is positive out to a hundred lags. For correlated data the effective
sample size is far below the nominal $n$, so KS $p$-values are severely
anti-conservative.
}

\workedarith{
In one dimension, $W_1$ is the average gap between sorted samples. Sort both, pair
by rank, and average the absolute differences.

\medskip
\textbf{A small example.} Real $\{1, 3, 5\}$, synthetic $\{2, 3, 4\}$:
\[
  W_1 = \tfrac13\bigl(|1-2| + |3-3| + |5-4|\bigr) = \tfrac13(1 + 0 + 1) = 0.667 .
\]

\medskip
\textbf{The shuffled control.} $W_1 = 0.000\text{e}{+}00$, exactly. The shuffled
series is a permutation of the real one, so sorting produces identical sequences
and every paired difference is zero.

\medskip
\textbf{A trained generator.} $W_1 \approx 0.001$--$0.005$.

\medskip
\textbf{A floor worth knowing about.} Suppose a generator collapses to emitting
the constant $\mu$, the real mean. Then $W_1$ is the mean absolute deviation of
the real data --- for BOVESPA, $0.00770$. That is the \emph{worst} a
constant-output model can score. A model producing genuine but imperfect returns
can easily exceed $0.00770$, so on this metric alone a collapsed model can beat a
working one. Chapter~\ref{ch:evaluation} returns to this.
}

\formaldef{
The \textbf{Wasserstein-$p$ distance} is
\[
  W_p(P,Q) = \left(\inf_{\gamma \in \Gamma(P,Q)}
    \E_{(X,Y)\sim\gamma}\norm{X-Y}^{p}\right)^{1/p} ,
\]
the infimum over couplings $\gamma$ whose marginals are $P$ and $Q$. In one
dimension with $p = 1$ this reduces to
\[
  W_1(P,Q) = \int_{-\infty}^{\infty}\bigl|F_P(x) - F_Q(x)\bigr|\,dx
           = \int_0^1 \bigl|F_P^{-1}(u) - F_Q^{-1}(u)\bigr|\,du .
\]
}

\analogybreak{
$W_1$ is permutation-invariant: $W_1(\text{real}, \text{shuffled}) = 0$ because
sorting erases order. That is simultaneously correct --- it genuinely measures
distributional proximity --- and a hard limitation, since it is blind to temporal
structure by construction. It belongs in the fidelity family for exactly that
reason, and no amount of care in computing it will make it detect a shuffle.
}

\whyhere{
\texttt{wasserstein} is one of the seven fidelity metrics. The shuffled-control
audit found it giving a perfect score to a series with no temporal structure at
all, which is what motivated splitting the composite into fidelity and temporal
families. The full optimal-transport treatment --- Kantorovich--Rubinstein duality,
and why $W_1$ is the right objective for a GAN --- belongs with the generative
models and is developed there.
}
